\documentclass[10pt, letterpaper]{article}

% Packages:
\usepackage[
    ignoreheadfoot, % set margins without considering header and footer
    top=0.5 cm, % seperation between body and page edge from the top
    bottom=0.5 cm, % seperation between body and page edge from the bottom
    left=1.8 cm, % seperation between body and page edge from the left
    right=1.8 cm, % seperation between body and page edge from the right
    footskip=1.0 cm, % seperation between body and footer
    % showframe % for debugging 
]{geometry} % for adjusting page geometry
\usepackage{titlesec} % for customizing section titles
\usepackage{tabularx} % for making tables with fixed width columns
\usepackage{array} % tabularx requires this
\usepackage[dvipsnames]{xcolor} % for coloring text
\definecolor{primaryColor}{RGB}{0, 0, 150} % define primary color
\usepackage{enumitem} % for customizing lists
\usepackage{fontawesome5} % for using icons
\usepackage{amsmath} % for math
\usepackage{natbib,bibentry} % for \bibentry
\usepackage[
    pdftitle={Sofya Dymchenko CV},
    pdfauthor={Sofya Dymchenko},
    pdfcreator={LaTeX with RenderCV},
    colorlinks=true,
    urlcolor=primaryColor
]{hyperref}
\usepackage[pscoord]{eso-pic} % for floating text on the page
\usepackage{calc} % for calculating lengths
\usepackage{bookmark} % for bookmarks
\usepackage{lastpage} % for getting the total number of pages
\usepackage{changepage} % for one column entries (adjustwidth environment)
\usepackage{paracol} % for two and three column entries
\usepackage{ifthen} % for conditional statements
\usepackage{needspace} % for avoiding page brake right after the section title
\usepackage{iftex} % check if engine is pdflatex, xetex or luatex

% % Ensure that generate pdf is machine readable/ATS parsable:
% \ifPDFTeX
%     \input{glyphtounicode}
%     \pdfgentounicode=1
%     \usepackage[T1]{fontenc}
%     \usepackage[utf8]{inputenc}
%     \usepackage{lmodern}
% \fi

\usepackage{charter}

% Some settings:
\raggedright
\AtBeginEnvironment{adjustwidth}{\partopsep0pt} % remove space before adjustwidth environment
\pagestyle{empty} % no header or footer
\setcounter{secnumdepth}{0} % no section numbering
\setlength{\parindent}{0pt} % no indentation
\setlength{\topskip}{0pt} % no top skip
\setlength{\columnsep}{0.15cm} % set column seperation
\pagenumbering{gobble} % no page numbering

\titleformat{\section}{\needspace{4\baselineskip}\bfseries\large}{}{0pt}{}[\vspace{1pt}\titlerule]

\titlespacing{\section}{
    % left space:
    -1pt
}{
    % top space:
    0.3 cm
}{
    % bottom space:
    0.2 cm
} % section title spacing

\renewcommand\labelitemi{$\vcenter{\hbox{\small$\bullet$}}$} % custom bullet points
\newenvironment{highlights}{
    \begin{itemize}[
        topsep=0.10 cm,
        parsep=0.10 cm,
        partopsep=0pt,
        itemsep=0pt,
        leftmargin=0 cm + 10pt
    ]
}{
    \end{itemize}
} % new environment for highlights


\newenvironment{highlightsforbulletentries}{
    \begin{itemize}[
        topsep=0.10 cm,
        parsep=0.10 cm,
        partopsep=0pt,
        itemsep=0pt,
        leftmargin=10pt
    ]
}{
    \end{itemize}
} % new environment for highlights for bullet entries

\newenvironment{onecolentry}{
    \begin{adjustwidth}{
        0 cm + 0.00001 cm
    }{
        0 cm + 0.00001 cm
    }
}{
    \end{adjustwidth}
} % new environment for one column entries

\newenvironment{twocolentry}[2][]{
    \onecolentry
    \def\secondColumn{#2}
    \setcolumnwidth{\fill, 4.5 cm}
    \begin{paracol}{2}
}{
    \switchcolumn \raggedleft \secondColumn
    \end{paracol}
    \endonecolentry
} % new environment for two column entries

\newenvironment{threecolentry}[3][]{
    \onecolentry
    \def\thirdColumn{#3}
    \setcolumnwidth{, \fill, 4.5 cm}
    \begin{paracol}{3}
    {\raggedright #2} \switchcolumn
}{
    \switchcolumn \raggedleft \thirdColumn
    \end{paracol}
    \endonecolentry
} % new environment for three column entries

\newenvironment{header}{
    \setlength{\topsep}{0pt}\par\kern\topsep\centering\linespread{1.5}
}{
    \par\kern\topsep
} % new environment for the header

\newcommand{\placelastupdatedtext}{% \placetextbox{<horizontal pos>}{<vertical pos>}{<stuff>}
  \AddToShipoutPictureFG*{% Add <stuff> to current page foreground
    \put(
        \LenToUnit{\paperwidth-2 cm-0 cm+0.05cm},
        \LenToUnit{\paperheight-1.0 cm}
    ){\vtop{{\null}\makebox[0pt][c]{
        \small\color{gray}\textit{Last updated in September 2024}\hspace{\widthof{Last updated in September 2024}}
    }}}%
  }%
}%

% save the original href command in a new command:
\let\hrefWithoutArrow\href

\begin{document}

    \nobibliography{my}
    \bibliographystyle{abbrvnat}

    \newcommand{\AND}{\unskip
        \cleaders\copy\ANDbox\hskip\wd\ANDbox
        \ignorespaces
    }
    \newsavebox\ANDbox
    \sbox\ANDbox{$|$}

    \begin{header}
        \fontsize{25 pt}{25 pt}\selectfont Sofya Dymchenko, Ph.D.

        \vspace{5 pt}

        \normalsize
        \mbox{Grenoble, France}%
        \kern 5.0 pt%
        \AND%
        \kern 5.0 pt%
        \mbox{\hrefWithoutArrow{mailto:sofya.dymchenko@gmail.com}{sofya.dymchenko@gmail.com}}%
        \kern 5.0 pt%
        \AND%
        \kern 5.0 pt%
        \mbox{\hrefWithoutArrow{https://svdcvt.github.io}{svdcvt.github.io}}%
    \end{header}

    \vspace{5 pt}% - 0.3 cm}

    % \section{Summary}
    % Applied researcher in \textbf{AI for Science and HPC}, specialising in data-driven surrogate modelling for scientific computing from data-centric and compute efficiency perspectives (\textbf{active learning, PINNs, distributed streamed training, foundation models}). 
    \textbf{Summary.} Applied Scientist with PhD in Computer Science, specialising in deep learning for scientific computing, data-centric machine learning, and high-performance computing (HPC), with publications at NeurIPS. Experience in engineering research software, designing and conducting extensive experiments, supervising small teams, and communicating with technical and non-technical audiences. 
    \section{Work Experience}

        \begin{twocolentry}{Mar 2022 -- Apr 2026}
            \textbf{Doctoral Researcher}, Grenoble, France
        \end{twocolentry}
        Datamove team, INRIA and University of Grenoble-Alpes (UGA)
        \vspace{0.10 cm}
        \begin{onecolentry}
            \begin{highlights}
                \item Pioneered a research line on query-synthetic active learning for deep PDE surrogates; initiated and led the development of the active sampling branch of \hrefWithoutArrow{https://melissa.gitlabpages.inria.fr/melissa/}{Melissa}, an open-source HPC/AI framework.
                \item Designed and implemented Breed [2, 3] for on-the-fly simulation data curation; achieved up to 10x lower surrogate errors under fixed compute budgets; orchestrated GPU/CPU experiments on Leonardo supercomputer.
                \item Extended this approach with generative modelling: conceived and co-developed scalable diffusion-based OGAS method [1] applied to conditional transformer-based surrogates and neural operators on 2D PDEs.
                \item Supervised two engineers and two MSc students on Melissa R\&D, including work on online pretraining of foundation models for PDE surrogates.
                % \item Supervised 2 engineers and 2 master's students contributing to the Melissa framework, including project on online pretraining of foundation models for PDE surrogates.
                % \item Introduced an active sampling method for PINNs training [3], which outperformed prior baselines by reducing training time by up tp 5x and lowering average errors by up to 2x for 1D Allen-Cahn PDE.
                % \item Contributed to an open-source HPC\&AI framework \hrefWithoutArrow{https://gitlab.inria.fr/melissa/melissa}{Melissa}; supervised 2 engineers and 2 master's students contributing to the project.
                % \item Developed active learning methods for deep PDE surrogates: an offline PINN training strategy that outperformed prior baselines [3], and a pioneering online framework that improves generalisation of a data-driven surrogate by curating data during training [1,2]. 
                % \item Publications: 2 first-author workshop papers (NeurIPS'23, SC'24), 1 under review.
            \end{highlights}
        \end{onecolentry}

        \vspace{0.1 cm}

        \begin{twocolentry}{Jun 2021 -- Dec 2021}
            \textbf{Research Intern}, Moscow, Russia
        \end{twocolentry}
        Computational Intelligence Lab, Skolkovo Institute of Science and Technology (Skoltech)
        \vspace{0.01 cm}
        \begin{onecolentry}
            \begin{highlights}
                \item Extended master's thesis work on optimal transport to computational geometry and computer vision.
            \end{highlights}
        \end{onecolentry}

        \vspace{0.05 cm}

        \begin{twocolentry}{Feb 2020 -- Aug 2020}
            \textbf{Research Intern}, Moscow, Russia
        \end{twocolentry}
        Media Algorithms Laboratory, Huawei R\&D Center
        \vspace{0.01 cm}
        \begin{onecolentry}
            \begin{highlights}
                \item Trained generative deep learning models with attention for personalized noise reduction in audio signals.
            \end{highlights}
        \end{onecolentry}

% \section{Teaching}

%     \begin{twocolentry}{Sep 2019 -- Feb 2024}
%         \textbf{Teaching Assistant \& Seminar Tutor}
%     \end{twocolentry}
%     \vspace{0.10 cm}
%     \begin{onecolentry}
%         \begin{highlights}
%             \item Led Python programming course for 3 consecutive years (100+ students): designed curriculum, managed teaching assistants, created assignments, and delivered lectures; awarded ``Best Course'' (2020).
%             \item Taught seminars in Data Science, Machine Learning, and Python for bachelor, master, and professional development students at HSE (2019-2021) and UGA (2022-2024).
%         \end{highlights}
%     \end{onecolentry}

    \section{Education}

        \begin{twocolentry}{Mar 2022 -- Apr 2026}
            \textbf{PhD in Computer Science}, Grenoble, France
        \end{twocolentry}
        Laboratory of Informatics of Grenoble (LIG), UGA
        \vspace{0.1 cm}
        \begin{onecolentry}
            \begin{highlights}
                \item Thesis: High Performance Online Deep Neural Network Training from Synthetic Data with Active Learning for Scientific Computing. The manuscript is rated ``Outstanding'' (top 5\%) by an external reviewer \href{https://svdcvt.github.io/phddefence/}{[link]}.
            \end{highlights}
        \end{onecolentry}
        \vspace{0.1 cm}

        \begin{twocolentry}{Sep 2019 -- Jun 2021}
            \textbf{MSc in Applied Mathematics and Computer Science}, Moscow, Russia
        \end{twocolentry}
        Statistical Learning Theory program, Skoltech jointly with HSE; with honors, GPA: 4.89/5
        \vspace{0.05 cm}
        \begin{onecolentry}
            \begin{highlights}
                \item Implemented a modified optimal transport algorithm, with complexity reduced from quadratic to log-linear, and its application to computational mesh generation
                 \hrefWithoutArrow{https://drive.google.com/drive/folders/17uK18UkpTXBUYaWh4q8fphDWdXxe_M9z?usp=drive_link}{[link]}.
                \item Managed team delivery of winning data science solutions for marketing applications, from initial concept to final presentation, in industry-supported university hackathons.
            \end{highlights}
        \end{onecolentry}

        \vspace{0.05 cm}

        \begin{twocolentry}{Sep 2015 -- Jun 2019}
            \textbf{BSc in Mathematics}, Moscow, Russia
        \end{twocolentry}
            National Research University Higher School of Economics (HSE), GPA: 7.87/10; Minor in Data Science
        % \vspace{0.05 cm}
        % \begin{onecolentry}
        %     \begin{highlights}
        %         \item Thesis: Collected and preprocessed 10-hours noisy YouTube dataset from scratch for rare languages multi-classification; trained multimodal CNN-RNN models achieving 77.6\% accuracy \hrefWithoutArrow{https://drive.google.com/file/d/1rkIUd2kjUYFmPsmkdUjKxfgaNsbIka4g/view?usp=sharing}{[link]}. 
        %     \end{highlights}
        % \end{onecolentry}


    \section{Additional Experience \& Skills}
        \begin{onecolentry}
            \textbf{Technical skills:} Python (PyTorch, JAX, Matplotlib, Pandas, Seaborn, Scikit-Learn, SciPy, NumPy, WandB), Git, Bash, SLURM, Linux/HPC environments,  \LaTeX, C{+\!+} (basic).
        \end{onecolentry}

        \vspace{0.09 cm}

        \begin{onecolentry}
            \textbf{Languages:} English (fluent), French (intermediate), Russian (native). 
        \end{onecolentry}

        \vspace{0.09 cm}

        \begin{onecolentry}
            \textbf{Teaching} (2019--2024): led Python course (3 years, 100+ students) incl. curriculum, TA supervision, lectures, won a ``Best Course'' award (2020); taught Data Science/ML seminars to non-CS students (HSE, UGA).
        \end{onecolentry}

        \vspace{0.09 cm}

        \begin{onecolentry}
            \textbf{Academic service}: on-site organisation and outreach for GAP workshops \hrefWithoutArrow{https://gap2024.sciencesconf.org}{[2024]}\hrefWithoutArrow{https://indico.math.cnrs.fr/event/15724/}{[2026]}; reviewer for WANT workshops at NeurIPS'23, ICML'24; member of UGA doctoral school student council and LIG laboratory synergy.
        \end{onecolentry}    

    \section{Publications}
        \begin{onecolentry}
            \begin{highlightsforbulletentries}
                \item[[1\hspace{-0.5em}]] \bibentry{cesarLearningWhereSimulateGenerative2026} \hrefWithoutArrow{https://arxiv.org/pdf/2606.09949}{[preprint]}
                \item[[2\hspace{-0.5em}]] \bibentry{dymchenkoMelissaDLBreedDataEfficient2024} \hrefWithoutArrow{https://hal.science/hal-04712480/}{[paper]} \hrefWithoutArrow{https://youtu.be/WUfa6SunwGc?si=zjwIf1VKYd00T9lH}{[video]}
                \item[[3\hspace{-0.5em}]] \bibentry{dymchenkoLossdrivenSamplingHardtolearn2023} \hrefWithoutArrow{https://hal.science/hal-04305233/}{[paper]}
            \end{highlightsforbulletentries}
        \end{onecolentry}
\end{document}